%【紙面設定】
%b4,横書き,2段
\documentclass[paper=b4j,landscape,twocolumn,fleqn]{jlreq}
%横書き用の調整
\special{papersize=\the\paperwidth,\the\paperheight}
%間に線をいれる
\setlength{\columnseprule}{0.4pt}
%ページ番号消す
\pagestyle{empty}
%ページ余白の調整
\usepackage[margin=10truemm]{geometry}
%画像挿入用のパッケージ
\usepackage[dvipdfmx]{graphicx}
%数学用のパッケージ
\usepackage{amsmath}
% 見出しの設定：行取りを1行にし、前後のスペースを極限まで削る
\ModifyHeading{section}{
  lines=1,
  before_space=10pt, % セクションの上の余白（適度に残すと見やすい）
  after_space=0pt    % セクションのしたの余白をゼロにする
}


%【本文】
\begin{document}
%1段落目
\section{方程式の計算(xを求めよ)}
\begin{align*}
    x + 5 &= 8 & 3x &= 12 & x + 10 &= 4 \\
    2x - 3 &= 7 & 15 - x &= 20 & 4x + 16 &= 0 \\
    3x + 4 &= x + 10 & 2x - 5 &= 5x + 7 & 4x - 8 &= 6x + 2 \\
    2(x - 3) &= 4 & 3(x + 4) &= -6 & 5(x + 1) &= 2x - 4 \\
    0.2x &= 4 & \frac{x}{3} - 2 &= -4 & 0.5x + 3 &= 1 
\end{align*}
\section{式の展開(式を展開せよ)}
\begin{align*}
    5(x+3) &= & (x+2)(x+6) &= \\
    -2(4x-7) &= & (x-3)(x+5) &= \\
    (x+8)^2 &= & (x-9)^2 &= \\
    (x+12)(x-12) &= & (2x+3)(x+2) &= \\
    (3x-1)(2x+5) &= & (4x-3)^2 &= 
\end{align*}
\section{因数分解(因数分解せよ)}
\begin{align*}
    2x+10 &= & x^2+7x+10 &= \\
    x^2-2x-8 &= & x^2-5x+4 &= \\
    x^2+6x+9 &= & x^2-25 &= \\
    x^2+x-20 &= & x^2+10x+25 &= \\
    x^2-81 &= & x^2+13x+42 &= 
\end{align*}
\section{文字と式(文章を元に式を作成せよ)}
\begin{align*}
　&1.　a円の商品を5個、b円の商品を3個買った。合計金額はc円である。\\
　&2.　1000円札を出して、x円の商品を3個買った。おつりはy円である。\\
　&3.　分速v\,\text{m}でt分間歩いた。進んだ道のりはd\,\text{m}である。\\
　&4.　100\,\text{L}入った水槽から、毎分x\,\text{L}ずつt分間水を抜いた。残りの水の量はy\,\text{L}である。\\
　&5.　x円の商品2個とy円の商品5個の合計金額に、10\%の消費税を加えるとz円になる。\\
　&6.　4回のテストの点数がそれぞれa, b, c, d点だった。平均点はm点である。\\
　&7.　現在地xから、負の方向に速さvでt秒間移動した。到達した地点の座標はpである。\\
　&8.　定価a円の商品をx\%引きで10個買った。合計金額はb円である。
\end{align*}

%2段落目
\newpage
\section{方程式の計算(xを求めよ)}
\begin{align*}
    x + 5 &= 8 & 3x &= 12 & x + 10 &= 4 \\
    2x - 3 &= 7 & 15 - x &= 20 & 4x + 16 &= 0 \\
    3x + 4 &= x + 10 & 2x - 5 &= 5x + 7 & 4x - 8 &= 6x + 2 \\
    2(x - 3) &= 4 & 3(x + 4) &= -6 & 5(x + 1) &= 2x - 4 \\
    0.2x &= 4 & \frac{x}{3} - 2 &= -4 & 0.5x + 3 &= 1 
\end{align*}
\section{式の展開(式を展開せよ)}
\begin{align*}
    5(x+3) &= & (x+2)(x+6) &= \\
    -2(4x-7) &= & (x-3)(x+5) &= \\
    (x+8)^2 &= & (x-9)^2 &= \\
    (x+12)(x-12) &= & (2x+3)(x+2) &= \\
    (3x-1)(2x+5) &= & (4x-3)^2 &= 
\end{align*}
\section{因数分解(因数分解せよ)}
\begin{align*}
    2x+10 &= & x^2+7x+10 &= \\
    x^2-2x-8 &= & x^2-5x+4 &= \\
    x^2+6x+9 &= & x^2-25 &= \\
    x^2+x-20 &= & x^2+10x+25 &= \\
    x^2-81 &= & x^2+13x+42 &= 
\end{align*}
\section{文字と式(文章を元に式を作成せよ)}
\begin{align*}
　&1.　a円の商品を5個、b円の商品を3個買った。合計金額はc円である。\\
　&2.　1000円札を出して、x円の商品を3個買った。おつりはy円である。\\
　&3.　分速v\,\text{m}でt分間歩いた。進んだ道のりはd\,\text{m}である。\\
　&4.　100\,\text{L}入った水槽から、毎分x\,\text{L}ずつt分間水を抜いた。残りの水の量はy\,\text{L}である。\\
　&5.　x円の商品2個とy円の商品5個の合計金額に、10\%の消費税を加えるとz円になる。\\
　&6.　4回のテストの点数がそれぞれa, b, c, d点だった。平均点はm点である。\\
　&7.　現在地xから、負の方向に速さvでt秒間移動した。到達した地点の座標はpである。\\
　&8.　定価a円の商品をx\%引きで10個買った。合計金額はb円である。
\end{align*}
\end{document}


